\documentclass[12pt]{article}
\usepackage[utf8]{inputenc}
\usepackage[T1]{fontenc}
\usepackage[margin=1in]{geometry}
\usepackage{setspace}
\usepackage{xcolor}
\usepackage[colorlinks=true, urlcolor=blue!60!black, linkcolor=blue!60!black]{hyperref}
\setlength{\parindent}{2em}
\setlength{\parskip}{1.25em}
\renewcommand{\baselinestretch}{1.0}
\title{\textit{Clinton v. City of New York}, 524 U.S. 417 (1998)}
\author{Case Brief}
\date{}

\begin{document}
\maketitle

\section{Case Metadata}
\begin{description}
  \item[Citation:] \textit{Clinton, President of the United States v. City of New York}, 524 U.S. 417 (1998) (No.~97-1374).
  \item[Argued / Decided:] Argued April 27, 1998; Decided June 25, 1998.
  \item[Vote:] 6--3 (affirming the U.S. District Court for the District of Columbia, 985 F. Supp. 168 (D.D.C. 1998)).
  \item[Author:] Justice John Paul Stevens (joined by Chief Justice Rehnquist and Justices Kennedy, Souter, Thomas, and Ginsburg).
  \item[Concurrences / Dissents:] Justice Kennedy (concurring, 524 U.S. at 449); Justice Scalia (concurring in part and dissenting in part, joined by Justice O'Connor, and joined by Justice Breyer as to Part III, 524 U.S. at 453); Justice Breyer (dissenting, joined by Justices O'Connor and Scalia as to Part III, 524 U.S. at 469).
  \item[Chapter / Topic:] Separation of Powers and the Executive --- Presentment Clause, Line-Item Veto, and the Power of the Purse (U.S. Const. art.~I, \S~7, cl.~2).
  \item[Primary Source URLs:] \href{https://supreme.justia.com/cases/federal/us/524/417/}{Justia U.S. Supreme Court Center} (\url{https://supreme.justia.com/cases/federal/us/524/417/}); \href{https://www.law.cornell.edu/supct/html/97-1374.ZO.html}{Cornell LII} (\url{https://www.law.cornell.edu/supct/html/97-1374.ZO.html}).
\end{description}

\section{Facts and Procedural History}
In April 1996, Congress enacted the Line Item Veto Act (Act), 2 U.S.C. \S~691 \textit{et seq.}, effective January 1, 1997, empowering the President to ``cancel in whole'' three categories of provisions within five calendar days (excluding Sundays) after signing a bill into law: (1) any dollar amount of discretionary budget authority; (2) any item of new direct spending; or (3) any limited tax benefit. To cancel an item, the President had to determine that cancellation would ``(i) reduce the Federal budget deficit; (ii) not impair any essential Government functions; and (iii) not harm the national interest,'' and transmit a special message to Congress. Cancellation prevented the provision ``from having legal force or effect'' (\S\S~691e(4)(B)--(C)) unless Congress passed a ``disapproval bill'' by majority vote of both Houses (which remained subject to presidential veto). A ``lockbox'' provision (\S~691c) devoted resulting savings to deficit reduction.

An initial suit by six Members of Congress was dismissed for lack of Article III standing in \textit{Raines v. Byrd}, 521 U.S. 811 (1997). Less than two months later, on August 11, 1997, President William J. Clinton exercised the Act to cancel:
\begin{enumerate}
  \item \textbf{Section 4722(c) of the Balanced Budget Act of 1997} (an item of new direct spending), which had waived the federal government's statutory right to recoup up to \$2.6 billion in Medicaid provider taxes collected by New York State; and
  \item \textbf{Section 968 of the Taxpayer Relief Act of 1997} (a limited tax benefit), which had amended 26 U.S.C. \S~1042 to let owners of food refining and processing facilities defer capital gains taxes when selling stock to eligible farmers' cooperatives.
\end{enumerate}

Two sets of plaintiffs filed separate suits in the U.S. District Court for the District of Columbia under the Act's expedited review provision (\S~692(a)(1)): (1) the City of New York, the New York City Health and Hospitals Corporation (NYCHHC), hospital associations, and health-care unions, who faced automatic retroactive tax assessments under New York law if HHS denied New York's waiver requests; and (2) Snake River Potato Growers, Inc., an Idaho farmers' cooperative formed to buy potato processing plants, and its vice-chairman, Mike Cranney, whose negotiations to purchase a processing facility ended after the cancellation. The District Court consolidated the cases, found Article III standing, and held the Act unconstitutional under both the Presentment Clause and general separation-of-powers principles. The Supreme Court expedited direct review under \S~692(b)--(c).

\section{Issue(s) Presented}
\begin{enumerate}
  \item \textbf{Jurisdiction and Standing:} Did institutional appellees qualify as ``individuals'' under the Act's expedited-review provision (\S~692(a)(1)), and did the New York and Snake River appellees suffer concrete, particularized Article III injuries sufficient to establish standing?
  \item \textbf{Constitutional Merits:} Do the cancellation procedures of the Line Item Veto Act violate the Presentment Clause, Article I, \S~7, cl.~2, by authorizing the President to unilaterally cancel portions of duly enacted statutes?
\end{enumerate}

\section{Holding}
Yes on jurisdiction and standing, and yes on the merits. Both sets of appellees had statutory and Article III standing, and the Line Item Veto Act's cancellation procedures violate the Presentment Clause because they authorize the President to unilaterally amend or repeal portions of duly enacted Acts of Congress without following the exclusive lawmaking procedures of Article I, \S~7.

\section{Reasoning}
Writing for a six-Justice majority, Justice Stevens first resolved statutory jurisdiction and Article III standing (Part III). Construing ``individual'' in \S~692(a)(1) as synonymous with ``person'' to avoid an absurd result, Stevens distinguished \textit{Raines v. Byrd} because these plaintiffs alleged concrete personal stakes rather than abstract institutional injuries: canceling \S~4722(c) immediately revived a multibillion-dollar contingent liability for New York and automatically assessed NYCHHC and New York City (analogous to an appellate court vacating a favorable judgment), while canceling \S~968 stripped Snake River of a targeted statutory ``bargaining chip'' in active acquisition negotiations (analogous to \textit{Bryant v. Yellen}, 447 U.S. 352 (1980)).

On the merits (Parts IV--VI), Stevens relied squarely on the formalist logic of \textit{INS v. Chadha}, 462 U.S. 919 (1983):
\begin{itemize}
  \item \textbf{Functional Equivalent of Partial Repeal:} Because cancellation prevented \S~4722(c) and \S~968 ``from having legal force or effect,'' the President had, in both legal and practical effect, amended two Acts of Congress by repealing portions of each. Under \textit{Chadha}, statutory repeal---no less than enactment---must conform to Article I, \S~7.
  \item \textbf{Constitutional Return vs. Statutory Cancellation:} Under Article I, \S~7, cl.~2, a constitutional ``return'' (veto) occurs \emph{before} a bill becomes law and applies to the \emph{entire} bill; statutory cancellation occurs \emph{after} a bill becomes law and strikes only a \emph{part}. Constitutional silence on unilateral presidential repeal is ``equivalent to an express prohibition,'' corroborated by George Washington's understanding that a President must ``approve all the parts of a Bill, or reject it in toto.''
  \item \textbf{Rejection of the ``Lockbox'' Defense:} That canceled items retained a budgetary scoring effect under the lockbox provision did not alter the fact that cancellation rendered both provisions completely inoperative as to appellees.

\item \textbf{Distinguishing \textit{Field v. Clark} (1892) and Traditional Spending Discretion:} Rejecting the Government's delegation defense, Stevens identified three critical differences between the Tariff Act upheld in \textit{Field v. Clark}, 143 U.S. 649 (1892), and the Line Item Veto Act: (1) \textit{Field}'s tariff suspension was contingent on post-enactment foreign trade conditions, whereas a five-day cancellation rested on the exact same facts Congress had just evaluated; (2) the President in \textit{Field} had a mandatory duty to suspend once the contingency arose, whereas the Act preserved unfettered presidential discretion; and (3) \textit{Field} executed congressional policy in foreign affairs, whereas cancellation rejected Congress's policy judgment. Likewise, traditional ``sums not exceeding'' appropriations never authorized the President ``to change the text of duly enacted statutes.''
  \item \textbf{Article V Exclusivity:} Allowing the President to produce ``Public Law 105-33 as modified by the President'' creates a law whose text was never voted on by either House or presented to the President. Altering that procedure requires a constitutional amendment under Article V.
\end{itemize}

\section{Concurring and Dissenting Opinions}
\begin{itemize}
  \item \textbf{Justice Kennedy (Concurring, 524 U.S. at 449--453):} Replied directly to Justice Breyer's claim that interbranch cooperation poses no threat to individual liberty. Kennedy emphasized that separation of powers protects liberty on both a horizontal axis (between branches) and a vertical axis (between government and citizens). Giving the President sole power to target specific states, taxpayers, or groups for punishment or leverage undermines political liberty, and ``one Congress cannot yield up its own powers, much less those of other Congresses to follow.''
  \item \textbf{Justice Scalia (Concurring in Part and Dissenting in Part, joined by Justice O'Connor, and joined by Justice Breyer as to Part III, 524 U.S. at 453--469):} Disagreed on statutory jurisdiction (arguing ``individual'' means natural person, though he would have granted certiorari before judgment under Rule 11) and argued Snake River lacked Article III standing because its negotiations were too speculative under \textit{Simon v. Eastern Ky. Welfare Rights Org.} and \textit{Allen v. Wright}. On the merits of New York's spending claim (Part III), Scalia argued the Act satisfied Article I, \S~7 because the Balanced Budget Act was duly enacted before any cancellation occurred. The real question was nondelegation, not Presentment---and from the First Congress's ``sums not exceeding'' appropriations and Jefferson's unspent gunboat funds to \textit{Train v. City of New York}, 420 U.S. 35 (1975), Congress had always been permitted to authorize the President to ``decline to spend'' appropriated sums. Calling that power a ``Line Item Veto,'' Scalia quipped, ``succeeded in faking out the Supreme Court.''
  \item \textbf{Justice Breyer (Dissenting, joined by Justices O'Connor and Scalia as to Part III, 524 U.S. at 469--497):} Agreed appellees had standing, but rejected the majority's literalist syllogism. In Part III, Breyer argued that when the President cancels an item pursuant to a preexisting statutory proviso, he does not ``repeal'' or ``amend'' a law; he \emph{executes} the statute according to its terms (analogous to exercising a power of appointment under a will or trust). In Part IV, Breyer applied a functionalist separation-of-powers and nondelegation analysis, arguing that given a \$1.5 trillion federal budget and 500-page omnibus bills, the Act neither encroached on Congress (which could exempt any bill by majority vote) nor violated the ``intelligible principle'' test for either spending items or limited tax benefits.
\end{itemize}

\section{Key Quotations}
\begin{quote}
``The constitutional return takes place before the bill becomes law; the statutory cancellation occurs after the bill becomes law. The constitutional return is of the entire bill; the statutory cancellation is of only a part.'' \\
--- Justice Stevens, 524 U.S. at 439.
\end{quote}

\begin{quote}
``There are powerful reasons for construing constitutional silence on this profoundly important issue as equivalent to an express prohibition.'' \\
--- Justice Stevens, 524 U.S. at 439.
\end{quote}

\begin{quote}
``If the Line Item Veto Act were valid, it would authorize the President to create a different law---one whose text was not voted on by either House of Congress or presented to the President for signature\ldots If there is to be a new procedure in which the President will play a different role in determining the final text of what may `become a law,' such change must come not by legislation but through the amendment procedures set forth in Article V of the Constitution.'' \\
--- Justice Stevens, 524 U.S. at 448--449.
\end{quote}

\begin{quote}
``Failure of political will does not justify unconstitutional remedies\ldots Liberty is always at stake when one or more of the branches seek to transgress the separation of powers\ldots Abdication of responsibility is not part of the constitutional design.'' \\
--- Justice Kennedy, concurring, 524 U.S. at 449, 450, 452.
\end{quote}

\begin{quote}
``The title of the Line Item Veto Act, which was perhaps designed to simplify for public comprehension, or perhaps merely to comply with the terms of a campaign pledge, has succeeded in faking out the Supreme Court.'' \\
--- Justice Scalia, concurring in part and dissenting in part, 524 U.S. at 469.
\end{quote}

\section{Doctrinal Significance and Modern Legacy (Through 2026)}
\textit{Clinton v. City of New York} stands as the mirror image of \textit{INS v. Chadha}: together, the two decisions bookend Article I, \S~7 formalism, prohibiting both Congress (\textit{Chadha}) and the President (\textit{Clinton}) from altering statutory law outside full bicameralism and presentment, even when both political branches enthusiastically agree to the institutional shortcut. \textit{Clinton} also remains a major Article III standing precedent for ``contingent liability'' and ``statutory bargaining chip'' injuries, as well as the bedrock authority confirming that Congress cannot delegate a unilateral line-item cancellation power without an Article V amendment.

Through 2025--2026, \textit{Clinton} has returned to the center of high-stakes constitutional litigation over the \textbf{power of the purse, executive impoundment, and ``pocket rescissions''}:
\begin{enumerate}
  \item \textbf{Impoundment and the Impoundment Control Act (ICA) of 1974:} In 2025, the second Trump administration revived claims of broad executive authority to freeze or withhold congressionally appropriated funds, arguing in part that the ICA unconstitutionally restricts Article II executive power. Challengers and constitutional scholars have invoked \textit{Clinton} (alongside \textit{Train v. City of New York}, 420 U.S. 35 (1975), and Justice Scalia's remark in \textit{Clinton}, 524 U.S. at 468, that \textit{Train} proved President Nixon's asserted constitutional right to impound wrong) to argue that if Congress cannot even \emph{grant} the President unilateral power to cancel enacted appropriations, the President certainly cannot claim an inherent Article II power to do so.
  \item \textbf{Pocket Rescissions and the 2025--2026 Emergency Docket:} In August 2025, the administration transmitted a special rescission message under the ICA seeking to withhold roughly \$4 billion in expiring foreign-aid appropriations with fewer than 45 days left in the fiscal year---a ``pocket rescission'' designed to let funds expire before Congress acted. On September 26, 2025, in \href{https://www.supremecourt.gov/opinions/24pdf/25a269_bp7c.pdf}{\textit{Department of State v. AIDS Vaccine Advocacy Coalition}, No. 25A269} (consolidated with \textit{Global Health Council v. Trump}), the Supreme Court voted 6--3 on its emergency docket to stay a district court injunction ordering obligation of the funds, concluding on a preliminary basis that the ICA precluded respondents' APA suit and citing foreign-affairs harms, over Justice Kagan's dissent (joined by Justices Sotomayor and Jackson). The funds expired four days later. A second pocket rescission of \$810 million in September 2026 drew a GAO determination of illegality and new state-led litigation. Commentators (see, e.g., \href{https://www.scotusblog.com/2025/10/shadow-overruling-its-not-just-for-humphreys-executor/}{SCOTUSblog, Oct. 13, 2025}; \href{https://harvardlawreview.org/blog/2025/11/making-sense-of-the-emergency-appropriations-decisions/}{Harvard Law Review Blog, Nov. 11, 2025}) have warned that if judicial enforcement of appropriations acts is foreclosed while funds expire, executive pocket rescissions risk achieving functionally what \textit{Clinton} forbade formally.
\end{enumerate}

\section{Pedagogical Critiques and Exam Hypotheticals}
\textbf{Pedagogical Critiques:}
\begin{itemize}
  \item \textit{Formalism vs. Substance (Scalia's Critique):} Does \textit{Clinton} elevate drafting magic words over constitutional substance? If Congress can constitutionally appropriate ``up to \$10 billion, in the President's discretion'' for a list of projects (which Stevens concedes at 524 U.S. at 446--447 is permissible), why does authorizing the President to ``cancel'' a line item violate Article I, \S~7? Students should evaluate whether the formal distinction between ``declining to spend'' and ``erasing statutory text'' meaningfully constrains executive power or merely forces Congress to draft permissive appropriations.
  \item \textit{Non-Severability and Tax vs. Spending:} Note that Stevens declined to reach severability (524 U.S. at 448 n.43), while Scalia's defense of spending discretion did not clearly extend to canceling \emph{tax} provisions (which he avoided via standing). In class, ask whether ``declining to collect a tax'' or ``revoking a tax exemption'' is more clearly legislative than ``declining to spend.''
\end{itemize}

\textbf{Exam Hypothetical:} To curb deficits without triggering \textit{Clinton}, Congress passes the Fiscal Restraint Act of 2026, providing that every future federal appropriations statute shall be deemed to appropriate ``sums not exceeding'' the listed amounts, and authorizing the President, within 10 days of signing any appropriations bill, to issue a ``non-expenditure directive'' permanently barring any agency from obligating designated line items whenever the President finds that non-expenditure will ``promote fiscal discipline.'' The statute expressly states that it ``does not repeal or amend any statutory text.'' President X signs an infrastructure bill and immediately issues a non-expenditure directive blocking \$500 million earmarked for a state port authority. The port authority sues. \textit{Analysis:} The Government will invoke Justice Scalia's partial dissent in \textit{Clinton} and historical ``sums not exceeding'' appropriations to argue that the Act delegates spending discretion rather than authorizing statutory repeal. The port authority will argue under Stevens's majority opinion that: (1) a blanket, across-the-board, 10-day post-enactment cancellation power based on the same facts Congress just evaluated is the \emph{functional equivalent} of a partial repeal violating Article I, \S~7 regardless of statutory labels; (2) unlike \textit{Field v. Clark}, the directive rejects Congress's immediate policy choice rather than executing a post-enactment contingency; and (3) under Kennedy's concurrence and the nondelegation doctrine, a vague ``fiscal discipline'' standard across the entire federal budget abdicates Congress's Article I power of the purse.

\textbf{``Opposite Side'' Contrast:} Contrast \textit{Clinton} with \textit{Field v. Clark}, 143 U.S. 649 (1892) (and \textit{J.W. Hampton, Jr., \& Co. v. United States}, 276 U.S. 394 (1928)). In \textit{Field}, the Court upheld Congress's delegation to the President of authority to suspend duty-free import status for specified commodities whenever foreign nations imposed ``reciprocally unequal and unreasonable'' tariffs. Whereas \textit{Field} represents constitutionally permissible \emph{contingent legislation}---where Congress itself chooses the policy trigger and the President executes it upon subsequent factual developments---\textit{Clinton} marks the outer formalist boundary where Congress attempts to hand the President a blank check to rewrite duly enacted domestic statutes based on his own immediate policy disagreements.

\end{document}
